\documentclass[11pt]{article}
\usepackage[margin=1in]{geometry}
\usepackage{amsmath,booktabs,multirow,tabularx,longtable,float,multicol,array,siunitx}
\usepackage[table]{xcolor}
\usepackage{caption}
\begin{document}
\section*{PHYS 1301 Handout 4: Data, units and tables}
Read the procedure before the lab session. Record every measurement with its unit and an estimate of its uncertainty, and keep your raw data even when a trial looks wrong. Read the procedure before the lab session. Record every measurement with its unit and an estimate of its uncertainty, and keep your raw data even when a trial looks wrong. 
\subsection*{Warm-up}
The rotation matrix and the linear system below are used in Problem 3.
\[ R = \begin{pmatrix} \cos\theta & -\sin\theta & 0 \\ \sin\theta & \cos\theta & 0 \\ 0 & 0 & 1 \end{pmatrix},
\qquad \left[\begin{array}{ccc|c} 2 & 1 & -1 & 8 \\ -3 & -1 & 2 & -11 \\ -2 & 1 & 2 & -3 \end{array}\right] \]
\begin{align*}
x(t) &= x_0 + v_0 t + \tfrac{1}{2} a t^2 \\
v(t) &= v_0 + a t \\
v^2 &= v_0^2 + 2 a (x - x_0)
\end{align*}
\[ f(x) = \begin{cases} 0 & x < 0 \\ x^2 & 0 \le x \le 1 \\ 1 & x > 1 \end{cases} \]
1. Which quantity is a vector?
\begin{multicols}{2}\begin{enumerate}\renewcommand{\labelenumi}{(\alph{enumi})}
\item mass \item speed \item displacement \item energy
\end{enumerate}\end{multicols}

\begin{table}[H]\centering
\caption{Cart on an incline: measured times}

\begin{tabular}{c S[table-format=1.2] S[table-format=1.2] S[table-format=1.2] S[table-format=1.3]}
\toprule
{\textbf{Trial}} & {\textbf{Mass (kg)}} & {\textbf{Time (s)}} & {\textbf{Distance (m)}} & {\textbf{Speed (m/s)}} \\
\midrule
1 & 0.30 & 1.27 & 1.50 & 1.181 \\
2 & 0.35 & 1.34 & 1.50 & 1.119 \\
3 & 0.40 & 1.41 & 1.50 & 1.064 \\
4 & 0.45 & 1.48 & 1.50 & 1.014 \\
5 & 0.50 & 1.55 & 1.50 & 0.968 \\
6 & 0.55 & 1.62 & 1.50 & 0.926 \\
7 & 0.60 & 1.69 & 1.50 & 0.888 \\
8 & 0.65 & 1.76 & 1.50 & 0.852 \\
9 & 0.70 & 1.83 & 1.50 & 0.820 \\
10 & 0.75 & 1.90 & 1.50 & 0.789 \\
\bottomrule
\end{tabular}
\end{table}

Read the procedure before the lab session. Record every measurement with its unit and an estimate of its uncertainty, and keep your raw data even when a trial looks wrong. 
\begin{table}[H]\centering
\caption{Physical constants used in this course}
\begin{tabular}{llll}
\toprule
\textbf{Quantity} & \textbf{Symbol} & \textbf{Value} & \textbf{Unit} \\
\midrule
Speed of light & $c$ & $2.998 \times 10^{8}$ & m/s \\
Gravitational constant & $G$ & $6.674 \times 10^{-11}$ & N m\textsuperscript{2}/kg\textsuperscript{2} \\
Planck constant & $h$ & $6.626 \times 10^{-34}$ & J s \\
Boltzmann constant & $k_\mathrm{B}$ & $1.381 \times 10^{-23}$ & J/K \\
Elementary charge & $e$ & $1.602 \times 10^{-19}$ & C \\
Avogadro constant & $N_\mathrm{A}$ & $6.022 \times 10^{23}$ & 1/mol \\
\bottomrule
\end{tabular}
\end{table}

Read the procedure before the lab session. Record every measurement with its unit and an estimate of its uncertainty, and keep your raw data even when a trial looks wrong. 
\begin{table}[H]\centering
\caption{Thermal properties of selected materials}
\begin{tabular}{lrrrr}
\toprule
\multirow{2}{*}{\textbf{Material}} & \multirow{2}{*}{\textbf{Density (kg/m\textsuperscript{3})}} & \multicolumn{2}{c}{\textbf{Specific heat (J/(kg K))}} & \multirow{2}{*}{\textbf{Melting point (°C)}} \\
\cmidrule(lr){3-4}
 & & \textbf{Solid} & \textbf{Liquid} & \\
\midrule
Aluminium & 2700 & 900 & 1180 & 660 \\
Copper & 8960 & 385 & 530 & 1085 \\
Lead & 11340 & 128 & 140 & 327 \\
Water (ice) & 917 & 2090 & 4186 & 0 \\
\bottomrule
\end{tabular}
\end{table}

\begin{table}[H]\centering
\caption{Lab report grading rubric}

\begin{tabularx}{\textwidth}{|l|X|X|X|}
\hline
\textbf{Criterion} & \textbf{Excellent (4)} & \textbf{Good (3)} & \textbf{Needs work (1)} \\
\hline
Hypothesis & Clear, testable and linked to the theory in the pre-lab. & Testable but not linked to theory. & Missing or not testable. \\ \hline
Data & All measurements recorded with units and uncertainties. & Units present; some uncertainties missing. & Units or values missing. \\ \hline
Analysis & Correct graph with best-fit line and slope uncertainty. & Graph correct; uncertainty not propagated. & No graph or wrong variables. \\ \hline
Conclusion & Answers the question and discusses sources of error. & Answers the question only. & Does not answer the question. \\ \hline
\end{tabularx}
\end{table}

Read the procedure before the lab session. Record every measurement with its unit and an estimate of its uncertainty, and keep your raw data even when a trial looks wrong. 
\begin{table}[H]\centering
\caption{Course schedule}

\begin{tabular}{lllll}
\textbf{Week} & \textbf{Dates} & \textbf{Topic} & \textbf{Reading} & \textbf{Due} \\
1 & Jan 13-17 & Units and measurement & Ch. 1 & - \\
2 & Jan 20-24 & Motion in one dimension & Ch. 2 & HW 1 \\
3 & Jan 27-31 & Vectors and projectiles & Ch. 3 & HW 2 \\
4 & Feb 3-7 & Newton's laws & Ch. 4 & Lab 1 \\
5 & Feb 10-14 & Friction and circular motion & Ch. 5 & HW 3 \\
6 & Feb 17-21 & Work and energy & Ch. 6 & Quiz 1 \\
7 & Feb 24-28 & Momentum & Ch. 7 & HW 4 \\
8 & Mar 3-7 & Midterm review & - & Midterm \\
\end{tabular}
\end{table}

Read the procedure before the lab session. Record every measurement with its unit and an estimate of its uncertainty, and keep your raw data even when a trial looks wrong. 
\begin{longtable}{rrrr}
\caption{Values of the trigonometric functions}\\
\toprule
\textbf{Angle (deg)} & \textbf{sin} & \textbf{cos} & \textbf{tan} \\
\midrule
\endfirsthead
\toprule
\textbf{Angle (deg)} & \textbf{sin} & \textbf{cos} & \textbf{tan} \\
\midrule
\endhead
\bottomrule
\endfoot
0 & 0.0000 & 1.0000 & 0.0000 \\
2 & 0.0349 & 0.9994 & 0.0349 \\
4 & 0.0698 & 0.9976 & 0.0699 \\
6 & 0.1045 & 0.9945 & 0.1051 \\
8 & 0.1392 & 0.9903 & 0.1405 \\
10 & 0.1736 & 0.9848 & 0.1763 \\
12 & 0.2079 & 0.9781 & 0.2126 \\
14 & 0.2419 & 0.9703 & 0.2493 \\
16 & 0.2756 & 0.9613 & 0.2867 \\
18 & 0.3090 & 0.9511 & 0.3249 \\
20 & 0.3420 & 0.9397 & 0.3640 \\
22 & 0.3746 & 0.9272 & 0.4040 \\
24 & 0.4067 & 0.9135 & 0.4452 \\
26 & 0.4384 & 0.8988 & 0.4877 \\
28 & 0.4695 & 0.8829 & 0.5317 \\
30 & 0.5000 & 0.8660 & 0.5774 \\
32 & 0.5299 & 0.8480 & 0.6249 \\
34 & 0.5592 & 0.8290 & 0.6745 \\
36 & 0.5878 & 0.8090 & 0.7265 \\
38 & 0.6157 & 0.7880 & 0.7813 \\
40 & 0.6428 & 0.7660 & 0.8391 \\
42 & 0.6691 & 0.7431 & 0.9004 \\
44 & 0.6947 & 0.7193 & 0.9657 \\
46 & 0.7193 & 0.6947 & 1.0355 \\
48 & 0.7431 & 0.6691 & 1.1106 \\
50 & 0.7660 & 0.6428 & 1.1918 \\
52 & 0.7880 & 0.6157 & 1.2799 \\
54 & 0.8090 & 0.5878 & 1.3764 \\
56 & 0.8290 & 0.5592 & 1.4826 \\
58 & 0.8480 & 0.5299 & 1.6003 \\
60 & 0.8660 & 0.5000 & 1.7321 \\
62 & 0.8829 & 0.4695 & 1.8807 \\
64 & 0.8988 & 0.4384 & 2.0503 \\
66 & 0.9135 & 0.4067 & 2.2460 \\
68 & 0.9272 & 0.3746 & 2.4751 \\
70 & 0.9397 & 0.3420 & 2.7475 \\
72 & 0.9511 & 0.3090 & 3.0777 \\
74 & 0.9613 & 0.2756 & 3.4874 \\
76 & 0.9703 & 0.2419 & 4.0108 \\
78 & 0.9781 & 0.2079 & 4.7046 \\
80 & 0.9848 & 0.1736 & 5.6713 \\
82 & 0.9903 & 0.1392 & 7.1154 \\
84 & 0.9945 & 0.1045 & 9.5144 \\
86 & 0.9976 & 0.0698 & 14.3007 \\
88 & 0.9994 & 0.0349 & 28.6363 \\
90 & 1.0000 & 0.0000 & undefined \\
\end{longtable}

Read the procedure before the lab session. Record every measurement with its unit and an estimate of its uncertainty, and keep your raw data even when a trial looks wrong. 
\begin{table}[H]\centering
\caption{Data for the inner planets}
\begin{tabular}{l|rrrr}
 & \textbf{Mercury} & \textbf{Venus} & \textbf{Earth} & \textbf{Mars} \\
\hline
\textbf{Mass (10\textsuperscript{24} kg)} & 0.330 & 4.87 & 5.97 & 0.642 \\
\textbf{Radius (km)} & 2440 & 6052 & 6371 & 3390 \\
\textbf{Orbital period (days)} & 88.0 & 224.7 & 365.2 & 687.0 \\
\textbf{Surface gravity (m/s\textsuperscript{2})} & 3.7 & 8.9 & 9.8 & 3.7 \\
\end{tabular}
\end{table}

Read the procedure before the lab session. Record every measurement with its unit and an estimate of its uncertainty, and keep your raw data even when a trial looks wrong. 

\par\medskip\noindent
\begin{tabular}{@{}ll@{}}
Instructor & Dr. Alex Rivera \\
Office & Science Building 2.104 \\
Office hours & Tuesday and Thursday, 2:00-3:30 pm \\
Email & arivera@example.edu \\
\end{tabular}
\par\medskip

Read the procedure before the lab session. Record every measurement with its unit and an estimate of its uncertainty, and keep your raw data even when a trial looks wrong. 
\begin{table}[H]\centering
\caption{Labs by unit}
\begin{tabular}{|l|l|l|}
\hline
\textbf{Unit} & \textbf{Topic} & \textbf{Lab activity} \\ \hline
\multirow{3}{*}{Mechanics} & Kinematics & Photogate timing of a falling ball \\ \cline{2-3}
 & Forces & Atwood machine \\ \cline{2-3}
 & Energy & Spring constant from oscillations \\ \hline
\multirow{2}{*}{Waves} & Sound & Resonance tube \\ \cline{2-3}
 & Light & Diffraction grating \\ \hline
\end{tabular}
\end{table}

\begin{table}[H]\centering
\caption{Quiz scores of 40 students}
\rowcolors{2}{gray!15}{white}

\begin{tabular}{lrrr}
\textbf{Score interval} & \textbf{Frequency} & \textbf{Relative frequency} & \textbf{Cumulative frequency} \\
50-59 & 3 & 0.075 & 3 \\
60-69 & 7 & 0.175 & 10 \\
70-79 & 14 & 0.350 & 24 \\
80-89 & 11 & 0.275 & 35 \\
90-100 & 5 & 0.125 & 40 \\
Total & 40 & 1.000 &  \\
\end{tabular}
\end{table}

\subsection*{Problems}
\begin{enumerate}
\item Plot distance against time squared for Table 1 and find the acceleration.
\item Use Table 2 to compute the gravitational force between two 1.0 kg masses 1.0 m apart.
\item Solve the linear system in the warm-up and apply $R$ with $\theta = 30^\circ$ to the solution.
\end{enumerate}
\end{document}
